\section{0/1 Knapsack}
\label{sec:dp-01-knapsack}

\problemstatement{Given $n$ items, where item $i$ has weight
$\textit{wt}[i]$ and value $\textit{val}[i]$, and a knapsack of capacity
$W$, choose a subset of items (each usable \textbf{at most once}) whose
total weight does not exceed $W$, maximizing total value.}

\begin{patternbox}
This is the chapter's namesake, and it generalizes
Section~\ref{sec:dp-subset-sum-target}'s boolean reachability table in
the most direct way possible: instead of asking ``\emph{can} we reach
exactly capacity $\textit{target}$,'' we ask ``what is the
\emph{maximum value} achievable using \emph{at most} capacity $W$'' --
replacing the logical OR combination with a \textbf{maximum}, and adding
each picked item's \emph{value} (not just consuming its weight) to the
running total. Contrast this directly with
the Greedy chapter's Fractional Knapsack problem: there,
items could be split, which let a single greedy ratio-sort solve it
optimally in $O(n \log n)$; here, items are all-or-nothing, which
destroys that greedy argument and is exactly why this problem needs a
genuine DP table.
\end{patternbox}

\subsection*{Deriving the recurrence}

Let $f(\textit{index}, \textit{capacity})$ be the maximum value
achievable using items $0..\textit{index}$, subject to a weight budget
of $\textit{capacity}$. At item $\textit{index}$: either exclude it
($f(\textit{index}-1, \textit{capacity})$), or include it (only if
$\textit{wt}[\textit{index}] \le \textit{capacity}$, gaining
$\textit{val}[\textit{index}]$ plus
$f(\textit{index}-1, \textit{capacity}-\textit{wt}[\textit{index}])$) --
take whichever is larger:
\begin{gather*}
  f(\textit{index}, \textit{capacity}) = \max\Big(
  f(\textit{index}{-}1, \textit{capacity}),\ \textit{take} \Big), \\
  \textit{take} = \textit{wt}[\textit{index}] \le \textit{capacity}\ ?\
  \textit{val}[\textit{index}] + f(\textit{index}{-}1,
  \textit{capacity}-\textit{wt}[\textit{index}]) : -\infty.
\end{gather*}
Base case: $f(0, \textit{capacity}) = \textit{val}[0]$ if
$\textit{wt}[0] \le \textit{capacity}$, else $0$.

\subsection*{Approach 1: Recursion}

\begin{lstlisting}
#include <bits/stdc++.h>
using namespace std;

int knapsackRecursive(int index, int capacity, vector<int>& wt, vector<int>& val) {
    if (index == 0) {
        return (wt[0] <= capacity) ? val[0] : 0;
    }

    int notPick = knapsackRecursive(index - 1, capacity, wt, val);
    int pick = INT_MIN;
    if (wt[index] <= capacity) {
        pick = val[index] + knapsackRecursive(index - 1, capacity - wt[index], wt, val);
    }
    return max(notPick, pick);
}

int main() {
    vector<int> wt = {1, 3, 4, 5};
    vector<int> val = {1, 4, 5, 7};
    int n = wt.size(), W = 7;
    cout << knapsackRecursive(n - 1, W, wt, val) << endl; // 9
    return 0;
}
\end{lstlisting}

\begin{complexitybox}[Approach 1 Complexity]
\textbf{Time.} $O(2^n)$.
\textbf{Space.} $O(n)$ recursion stack.
\end{complexitybox}

\subsection*{Approach 2: Memoization}

\begin{lstlisting}
#include <bits/stdc++.h>
using namespace std;

int knapsackMemo(int index, int capacity, vector<int>& wt, vector<int>& val,
                  vector<vector<int>>& memo) {
    if (index == 0) {
        return (wt[0] <= capacity) ? val[0] : 0;
    }
    if (memo[index][capacity] != -1) return memo[index][capacity];

    int notPick = knapsackMemo(index - 1, capacity, wt, val, memo);
    int pick = INT_MIN;
    if (wt[index] <= capacity) {
        pick = val[index] + knapsackMemo(index - 1, capacity - wt[index], wt, val, memo);
    }
    return memo[index][capacity] = max(notPick, pick);
}

int main() {
    vector<int> wt = {1, 3, 4, 5};
    vector<int> val = {1, 4, 5, 7};
    int n = wt.size(), W = 7;
    vector<vector<int>> memo(n, vector<int>(W + 1, -1));
    cout << knapsackMemo(n - 1, W, wt, val, memo) << endl; // 9
    return 0;
}
\end{lstlisting}

\begin{complexitybox}[Approach 2 Complexity]
\textbf{Time.} $O(n \cdot W)$.
\textbf{Space.} $O(n \cdot W)$ memo $+$ $O(n)$ recursion stack.
\end{complexitybox}

\subsection*{Approach 3: Tabulation}

\begin{lstlisting}
#include <bits/stdc++.h>
using namespace std;

int knapsackTabulation(vector<int>& wt, vector<int>& val, int W) {
    int n = wt.size();
    vector<vector<int>> dp(n, vector<int>(W + 1, 0));

    for (int capacity = wt[0]; capacity <= W; capacity++) dp[0][capacity] = val[0];

    for (int index = 1; index < n; index++) {
        for (int capacity = 0; capacity <= W; capacity++) {
            int notPick = dp[index - 1][capacity];
            int pick = INT_MIN;
            if (wt[index] <= capacity) {
                pick = val[index] + dp[index - 1][capacity - wt[index]];
            }
            dp[index][capacity] = max(notPick, pick);
        }
    }
    return dp[n - 1][W];
}

int main() {
    vector<int> wt = {1, 3, 4, 5};
    vector<int> val = {1, 4, 5, 7};
    cout << knapsackTabulation(wt, val, 7) << endl; // 9
    return 0;
}
\end{lstlisting}

\subsection*{Dry run of the tabulation table}

\dryrun{Take $\textit{wt}=[1,3,4,5]$, $\textit{val}=[1,4,5,7]$, $W=7$.}

\begin{center}
\begin{tabular}{c|cccccccc}
$\textit{index}\backslash W$ & 0 & 1 & 2 & 3 & 4 & 5 & 6 & 7 \\ \hline
$0$ (wt$1$,val$1$) & 0 & 1 & 1 & 1 & 1 & 1 & 1 & 1 \\
$1$ (wt$3$,val$4$) & 0 & 1 & 1 & 4 & 5 & 5 & 5 & 5 \\
$2$ (wt$4$,val$5$) & 0 & 1 & 1 & 4 & 5 & 6 & 6 & 9 \\
$3$ (wt$5$,val$7$) & 0 & 1 & 1 & 4 & 5 & 7 & 8 & 9
\end{tabular}
\end{center}

At $\textit{index}=2$, $W=7$: not-pick gives $\textit{dp}[1][7]=5$; pick
gives $\textit{val}[2]+\textit{dp}[1][3] = 5+4=9$. Max is $9$. The final
answer is $\textit{dp}[3][7]=9$, achieved by items $1$ (wt$3$, val$4$)
and $2$ (wt$4$, val$5$) together: total weight $3+4=7\le7$, total value
$4+5=9$, matching the table.

\begin{complexitybox}[Approach 3 Complexity]
\textbf{Time.} $O(n \cdot W)$.
\textbf{Space.} $O(n \cdot W)$ table, no recursion stack.
\end{complexitybox}

\subsection*{Approach 4: Space Optimization}

\begin{lstlisting}
#include <bits/stdc++.h>
using namespace std;

int knapsackSpaceOptimized(vector<int>& wt, vector<int>& val, int W) {
    int n = wt.size();
    vector<int> prev(W + 1, 0);

    for (int capacity = wt[0]; capacity <= W; capacity++) prev[capacity] = val[0];

    for (int index = 1; index < n; index++) {
        vector<int> current(W + 1, 0);
        for (int capacity = 0; capacity <= W; capacity++) {
            int notPick = prev[capacity];
            int pick = INT_MIN;
            if (wt[index] <= capacity) {
                pick = val[index] + prev[capacity - wt[index]];
            }
            current[capacity] = max(notPick, pick);
        }
        prev = current;
    }
    return prev[W];
}

int main() {
    vector<int> wt = {1, 3, 4, 5};
    vector<int> val = {1, 4, 5, 7};
    cout << knapsackSpaceOptimized(wt, val, 7) << endl; // 9
    return 0;
}
\end{lstlisting}

\begin{complexitybox}[Approach 4 Complexity]
\textbf{Time.} $O(n \cdot W)$.
\textbf{Space.} $O(W)$ for the rolling row.
\end{complexitybox}

\begin{takeawaybox}
0/1 Knapsack is the most general member of this chapter's 0/1 family:
Subset Sum is the special case where every item's weight equals its
value (so reachability and value-maximization coincide), and Count
Subsets is the special case that counts qualifying selections instead of
optimizing a value. Recognizing 0/1 Knapsack as the ``parent'' recurrence
of everything from Section~\ref{sec:dp-subset-sum-target} onward
consolidates this entire half of the chapter into one shape.
\end{takeawaybox}
